\section{Reaktionskrafte}
Erst ist die maßstäbliche Abbildung des Konstruktions zu sehen:
\begin{figure}[H]
  \centering
  \begin{tikzpicture}[thick]
    \def\mscale{3};
    \def\a{0.5 * \mscale};
    \def\b{2.0 * \mscale};
    \def\c{2.0 * \mscale};
    \def\h{0.433 * \mscale};

    \def\kNscale{0.5};
    \def\G{3 * \kNscale};

    \def\kNmscale{1.25};
    \def\p{0.5 * \kNmscale};

    \coordinate (A) at (-\c, -\h);
    \coordinate (B) at (\a / 2, -\h);
    \coordinate (C) at (0, 0);
    \coordinate (D) at (\a, 0);
    \coordinate (E) at (\a + \b, 0);

    \draw[vec={red}] (E) -- ++(0, -\G) node[left] {$G$};

    \fill[pattern=vertical lines, pattern color=red] (C) rectangle ($(E) + (0, \p)$);
    \draw[red] (C) rectangle ($(E) + (0, \p)$);
    \draw[force={red}] ($(E) + (0, \p)$) -- (E) node[midway, right] {$p$};

    \draw (A) -- (C);
    \draw (D) -- (C);
    \draw (B) -- (C);
    \draw (D) -- (B);
    \draw (D) -- (E);
    \draw [dashed, thin, color=gray] (A) -- (B);
    \supporthinge{A}{0}
    \supporthinge{B}{0}
    \hinge{C}
    \draw (D) ++(0.175, 0) arc (0:180:0.175);
    \hinge{D}
  \end{tikzpicture}
  \caption{Maßstäblich korrekte Abbildung des Konstruktions}
\end{figure}
Um die Reaktionen im Punkt A und B zu bestimmen, kann die Konstruktion als ein Starrkörper angenommen werden:
\begin{figure}[H]
  \centering
  \begin{tikzpicture}[thick]
    \def\mscale{3};
    \def\a{0.5 * \mscale};
    \def\b{2.0 * \mscale};
    \def\c{2.0 * \mscale};
    \def\h{0.433 * \mscale};

    \def\kNscale{0.5};
    \def\G{3 * \kNscale};

    \def\kNmscale{1.25};
    \def\p{0.5 * \kNmscale};

    \draw[force={red}] (E) -- ++(0, -\G) node[left] {$G$};

    \fill[pattern=vertical lines, pattern color=red] (C) rectangle ($(E) + (0, \p)$);
    \draw[red] (C) rectangle ($(E) + (0, \p)$);
    \draw[force={red}] ($(E) + (0, \p)$) -- (E) node[midway, right] {$p$};

    \draw (A) -- (C);
    \draw (D) -- (C);
    \draw (B) -- (C);
    \draw (D) -- (B);
    \draw (D) -- (E);

    \draw[force={red}, <-] (A) -- ++(-1, 0) node[midway, above] {$A_X$};
    \draw[force={red}, <-] (A) -- ++(0, -1) node[left] {$A_Y$};
    \draw[force={red}, <-] (B) -- ++(-1, 0) node[midway, above] {$B_X$};
    \draw[force={red}, <-] (B) -- ++(0, -1) node[left] {$B_Y$};

  \end{tikzpicture}
  \caption{Abbildung der Reaktionskräfte}
\end{figure}

Man bekommt 3 Gleichgewichtsgleichungen, und 1 Gleichung weil Stab 5 nur Kräfte nur im Stabrichtung auswirken kann.

\begin{align}
  &\sum F_X = A_X + B_X = 0 \\
  &\sum F_Y = A_Y + B_Y - G - p \cdot (a + b) = 0 \\
  &\sum M_A = B_Y \cdot (c + \frac{a}{2}) - G \cdot (a + b + c) - p \cdot (a + b) \cdot (c + \frac{a + b}{2}) = 0\\
  &\frac{A_X}{A_Y} = \frac{c}{\frac{a}{2} \cdot tg(60^\circ)}
\end{align}

Daraus bekommt man die folgende Werte für die Reaktionskräfte:
\begin{align*}
  \boxed{
    \begin{aligned}
      A_X &= -16{,}422 \, \text{kN} \\
      A_Y &= -3{,}556 \, \text{kN}
  \end{aligned}}
  &\qquad
  \boxed{
    \begin{aligned}
      B_X &= 16{,}422 \, \text{kN} \\
      B_Y &= 7{,}806 \, \text{kN}
  \end{aligned}}
\end{align*}
